\documentclass[11pt]{article}
\usepackage[margin=1in]{geometry}
\usepackage{times}
\usepackage{hyperref}
\hypersetup{colorlinks=true, linkcolor=black, urlcolor=blue, citecolor=black}
\title{A Resonant Cavity Is Not a New Category: The Nassikas Thruster and the EmDrive Family Resemblance}
\author{Flyxion \\ \small Independent Researcher}
\date{}
\begin{document}
\maketitle

\begin{abstract}
The Nassikas thruster, proposed by John Nassikas and collaborators, claims that a superconducting resonant electromagnetic cavity can produce net thrust through an asymmetric interaction with a magnetic field, framed by its proponents as distinct from and more theoretically grounded than the EmDrive. This paper argues that despite differences in claimed mechanism, the device shares the EmDrive's core evidentiary weakness, a reported thrust at or near the noise floor of the measuring apparatus, no independent replication under improved isolation, and a theoretical justification that, on close reading, reduces to the same kind of appeal to an asymmetric force within a closed system that momentum conservation forbids regardless of the specific field configuration invoked to produce it.
\end{abstract}

\section{Introduction}

The Nassikas thruster proposes using a superconducting resonant cavity through which a magnetic field is passed, claiming that the interaction between a time-varying magnetic flux and the cavity's superconducting boundary produces a net force distinct in origin from the microwave radiation pressure asymmetry claimed by Shawyer's original EmDrive, and its proponents have specifically framed it as theoretically cleaner and less speculative, appealing to established superconductor electrodynamics rather than the more informally justified microwave cavity mechanics of the original EmDrive proposal. Reported thrust measurements remain in the micronewton range, comparable to and in some cases smaller than the EmDrive measurements addressed elsewhere in this series.

\section{The Family Resemblance the Framing Obscures}

Whatever specific electromagnetic interaction is invoked, superconducting boundary condition, microwave radiation pressure gradient, or any other asymmetric field effect, a closed system with no propellant expelled and no external force applied cannot experience a net change in its center of mass momentum, a consequence of momentum conservation that does not depend on which internal mechanism is proposed to evade it. Reframing the internal mechanism from microwave cavity asymmetry to superconducting flux interaction changes the physics being invoked without changing the conservation-law objection that applies regardless, since the objection concerns the system's boundary condition, closed, no propellant, no external force, rather than the specific internal process claimed to produce the resulting force.

\section{Why the Superconductor Framing Sounds More Credible Without Being More Supported}

Superconductor electrodynamics is a genuinely rigorous and well-established field, and invoking it lends a device an air of theoretical seriousness that a more informally justified proposal lacks, but rigor in the underlying physics used to describe a claimed interaction does not transfer to rigor in the claim that the interaction produces net external thrust unless the specific derivation is shown, step by step, to conserve momentum internally while still producing external motion, a demonstration the published Nassikas thruster literature has not supplied in a form that has satisfied independent theoretical reviewers, several of whom have identified the same category of error common across reactionless-drive proposals: treating an internal force pair as though only one member of the pair need be accounted for.

\section{The Measurement Problem Repeats}

As with the EmDrive, reported thrust levels for the Nassikas thruster sit close enough to the sensitivity floor of torsion-balance and similar thrust-measurement apparatus that thermal expansion during power-up, electromagnetic interference with the balance's own sensing elements, and outgassing in imperfect vacuum conditions remain live, undistinguished alternative explanations, and no independent laboratory outside the group originally proposing the device has published a replication attempt with improved isolation of the kind that resolved much of the ambiguity in the EmDrive literature by driving the measured effect toward zero.

\section{The Entropic Gravity Framing}

The Nassikas thruster's claimed mechanism is electromagnetic rather than gravitational, so the entropic account of gravitation does not directly bear on its central claim, but it does, as with the EmDrive, remove the common secondary appeal to the device pushing against "spacetime" or the "quantum vacuum" as a substrate with locally available momentum to borrow from, since on this account no such locally available substrate exists at any scale a laboratory cavity could access.

\section{Objections}

Proponents note correctly that the specific theoretical derivation differs meaningfully from Shawyer's original EmDrive proposal and should be evaluated independently rather than dismissed by association. This is a fair methodological point in principle, and the theoretical critique offered above engages the Nassikas-specific mechanism directly rather than relying on guilt by association; the conclusion that the derivation does not close the momentum-conservation gap it needs to close is a claim about this specific proposal's own published reasoning, not merely an inference from its resemblance to an earlier, differently justified device.

\section{Conclusion}

The Nassikas thruster's more rigorous-sounding theoretical vocabulary does not supply a demonstrated resolution to the same momentum-conservation objection that applies to every closed-system reactionless thrust claim, and its measured thrust remains, as with the EmDrive, within the range that improved experimental isolation has repeatedly shown to be explicable by mundane systematic effects rather than genuine anomalous force.

\end{document}
